\documentclass[11pt]{article}

\usepackage[margin=1.1in]{geometry}
\usepackage[T1]{fontenc}
\usepackage{amsmath,amssymb,amsthm,mathtools}
\usepackage{booktabs,array}
\usepackage{hyperref}
\usepackage{parskip}
\usepackage{microtype}
\usepackage{xcolor}
\usepackage{enumitem}
\usepackage{listings}
\usepackage{times}

\hypersetup{colorlinks=true, linkcolor=blue!60!black, citecolor=blue!60!black, urlcolor=blue!60!black}
\setlist[enumerate]{itemsep=0.3em, topsep=0.4em}
\setlist[itemize]{itemsep=0.3em, topsep=0.4em}

\lstset{
  language=Matlab,
  basicstyle=\ttfamily\small,
  columns=fullflexible,
  breaklines=true,
  keepspaces=true,
  showstringspaces=false,
  frame=single,
  framerule=0.3pt,
  rulecolor=\color{black!35},
  xleftmargin=0.6em,
  xrightmargin=0.6em,
  aboveskip=0.8em,
  belowskip=0.8em
}

\newtheorem{proposition}{Proposition}
\newtheorem{remark}{Remark}
\DeclareMathOperator*{\argmin}{arg\,min}
\newcommand{\E}{\mathbb{E}}
\newcommand{\Var}{\operatorname{Var}}
\newcommand{\Cov}{\operatorname{Cov}}
\newcommand{\R}{\mathbb{R}}

\title{\textbf{Inference for a Simulated Spatial Diffusion Estimator}\\[0.35em]
\large Analytical Sandwich Variance and Bootstrap Approaches}
\author{}
\date{}

\begin{document}
\maketitle
\tableofcontents
\bigskip

\section{Problem setup and estimator}

Consider a spatial diffusion process observed at $L$ sites. Each site
$\ell = 1, \ldots, L$ has known coordinates $s_\ell$ (or a mapping to grid
cells) and an observed arrival time $Y_\ell$.

Let $\theta \in \R^K$ denote the parameter vector of a stochastic diffusion
simulator. For a given $\theta$, the simulator produces a simulated arrival
(or adoption) time at site $\ell$ on simulation draw $n$:
\[
\hat Y_{\ell,n}(\theta), \qquad n=1,\ldots,N.
\]
Define the Monte Carlo (MC) average
\begin{equation}
\bar Y_{\ell,N}(\theta)
= \frac{1}{N}\sum_{n=1}^N \hat Y_{\ell,n}(\theta).
\label{eq:mcavg}
\end{equation}
The simulated nonlinear least squares (SNLS) estimator minimizes the mean
squared error between observed and simulated arrival times:
\begin{equation}
\hat\theta_{L,N}
= \argmin_{\theta \in \R^K}
Q_{L,N}(\theta),
\qquad
Q_{L,N}(\theta)
= \frac{1}{L}\sum_{\ell=1}^L \bigl(Y_\ell - \bar Y_{\ell,N}(\theta)\bigr)^2.
\label{eq:criterion}
\end{equation}
This is also a special case of simulated method of moments (SMM) with an
identity weighting matrix.

It is often helpful to distinguish the large-$N$ ``oracle'' target
\[
\mu_\ell(\theta) := \E\bigl[\hat Y_{\ell,1}(\theta)\bigr]
\]
from the finite-$N$ approximation $\bar Y_{\ell,N}(\theta)$. Define oracle
(respectively, feasible) residuals
\[
 e_\ell(\theta) := Y_\ell - \mu_\ell(\theta),
\qquad
 \hat e_\ell(\theta) := Y_\ell - \bar Y_{\ell,N}(\theta),
\]
and the oracle criterion
\[
\tilde Q_L(\theta) = \frac{1}{L}\sum_{\ell=1}^L e_\ell(\theta)^2.
\]
The oracle estimator is $\tilde\theta_L = \argmin_{\theta} \tilde Q_L(\theta)$.

\begin{remark}[Two inferential targets]
There are two natural inferential targets:
\begin{itemize}
  \item \textbf{Oracle (large-$N$) inference:} approximate the sampling distribution
        of $\tilde\theta_L$ (as if $\mu_\ell(\theta)$ were known).
  \item \textbf{Finite-$N$ inference:} approximate the sampling distribution of the
        reported estimator $\hat\theta_{L,N}$, which includes MC error.
\end{itemize}
In many applications one chooses $N$ large enough that MC error is secondary
and focuses on the oracle target.
\end{remark}

\section{Analytical variance: asymptotic normality and sandwich formula}

\subsection{Jacobian and curvature}

Let
\[
J_\ell(\theta) := \frac{\partial \mu_\ell(\theta)}{\partial \theta'} \in \R^{1\times K},
\qquad
J(\theta) := \begin{bmatrix}J_1(\theta)\\ \vdots \\ J_L(\theta)\end{bmatrix} \in \R^{L\times K}.
\]
A natural sample curvature matrix is the Gauss--Newton matrix
\begin{equation}
\hat\Gamma
= \frac{1}{L} J(\hat\theta)'J(\hat\theta) \in \R^{K\times K}.
\label{eq:gammahat}
\end{equation}
Its population counterpart is $\Gamma = \E[J_\ell(\theta_0)'J_\ell(\theta_0)]$.

\subsection{Asymptotic variance}

Define the oracle score contribution
\[
\psi_\ell(\theta_0)
:= -2 e_\ell(\theta_0)\, J_\ell(\theta_0)' \in \R^K.
\]
Under standard regularity conditions for nonlinear least squares
(identification, local smoothness of $\mu_\ell(\theta)$, nonsingularity of
$\Gamma$, and either site independence or an appropriate weak-dependence/spatial
mixing condition),
\begin{equation}
\sqrt{L}\,(\tilde\theta_L - \theta_0)
\xrightarrow{d}
\mathcal{N}(0, V),
\qquad
V = \Gamma^{-1} S\, \Gamma^{-1},
\label{eq:asymnormal}
\end{equation}
where $S$ is the long-run variance of the score contributions.

\paragraph{Site-i.i.d. (heteroskedastic) benchmark.}
If sites are independent (but not necessarily homoskedastic), then
\[
S = \E\bigl[ e_\ell(\theta_0)^2\,J_\ell(\theta_0)'J_\ell(\theta_0)\bigr],
\]
and the plug-in estimator is
\begin{equation}
\hat S_{\mathrm{iid}}
= \frac{1}{L} J(\hat\theta)' \operatorname{diag}(\hat e^2)\, J(\hat\theta).
\label{eq:meat_iid}
\end{equation}

\paragraph{Homoskedastic shortcut.}
If in addition $\E[e_\ell(\theta_0)^2 \mid X_\ell] = \sigma^2$ (strong), then
$S=\sigma^2\Gamma$ and $V=\sigma^2\Gamma^{-1}$. A simple estimate is
\begin{equation}
\hat\sigma^2 = \frac{1}{L-K}\sum_{\ell=1}^L \hat e_\ell^2.
\label{eq:sigma2hat}
\end{equation}

\paragraph{General sandwich variance.}
Given any consistent estimator $\hat S$ of $S$, define
\begin{equation}
\widehat{\Var}(\tilde\theta_L)
\approx \frac{1}{L}\, \hat\Gamma^{-1} \hat S\, \hat\Gamma^{-1}.
\label{eq:varhat_general}
\end{equation}
The matrix $\hat\Gamma^{-1} \hat S\, \hat\Gamma^{-1}$ estimates the asymptotic
variance of $\sqrt{L}(\tilde\theta_L-\theta_0)$; dividing by $L$ gives the
variance of $\tilde\theta_L$.

\subsection{Equivalent stacked notation}

Some presentations write the ``meat'' in terms of a stacked covariance matrix.
Let $g(\theta_0)=(g_1,\ldots,g_L)'$ where $g_\ell = Y_\ell-\mu_\ell(\theta_0)$.
Then
\[
\Omega := \Var\bigl(g(\theta_0)\bigr) \in \R^{L\times L}.
\]
With $J(\theta_0)$ defined as above, the asymptotic variance can be written as
\[
V = \Gamma^{-1}\, \Bigl(\lim_{L\to\infty} \tfrac{1}{L}J(\theta_0)'\Omega J(\theta_0)\Bigr)\, \Gamma^{-1}.
\]
Under i.i.d. homoskedastic residuals, $\Omega=\sigma^2 I_L$ and the expression
reduces to the homoskedastic shortcut.

\subsection{Simulation noise and the role of $N$}

At the true parameter, the feasible residual decomposes as
\[
\hat e_\ell(\theta_0)
= \underbrace{\bigl(Y_\ell - \mu_\ell(\theta_0)\bigr)}_{\text{data / model error}}
\;+
\underbrace{\bigl(\mu_\ell(\theta_0) - \bar Y_{\ell,N}(\theta_0)\bigr)}_{\text{MC error}}.
\]
The MC error term has variance of order $1/N$.

\begin{proposition}[Large-$N$ regime]
If $L/N \to 0$ (equivalently $\sqrt{L/N}\to 0$), then MC error is asymptotically
negligible relative to sampling error and
\[
\sqrt{L}(\hat\theta_{L,N}-\tilde\theta_L)=o_p(1).
\]
Consequently, $\hat\theta_{L,N}$ and $\tilde\theta_L$ share the same first-order
limit distribution in \eqref{eq:asymnormal}.
\end{proposition}

\begin{remark}[Practical interpretation]
In applications, $N$ is a computational tuning parameter.
\begin{itemize}
  \item If inference is meant to reflect uncertainty from the data (not from MC
        approximation), choose $N$ large enough that MC error is small.
  \item If one intentionally treats MC noise as part of the estimator, keep $N$
        fixed and interpret inference as describing the finite-$N$ estimator.
\end{itemize}
\end{remark}

\subsection{Spatial dependence across sites}

Treating sites as independent draws can be inappropriate in a diffusion
application. Spatial dependence can enter through shared measurement error,
shared omitted variables, or misspecification that varies smoothly in space.
If dependence is material, then the i.i.d. meat \eqref{eq:meat_iid} is not
appropriate.

Two practical dependence-robust alternatives are:
\begin{itemize}
  \item \textbf{Cluster / block robust sandwich.} Partition sites into spatial
        clusters (grid blocks, regions, or distance-based clusters) and use a
        cluster-robust meat.
  \item \textbf{Spatial HAC (Conley-type).} Estimate $S$ by weighting
        cross-products of score contributions by a distance kernel.
\end{itemize}
A cluster/block bootstrap (Section~\ref{sec:bootstrap}) is often the simplest
way to obtain dependence-robust inference.

\section{Bootstrap inference}
\label{sec:bootstrap}

Bootstrap inference re-estimates the model on resampled data. It can avoid
explicit Jacobian estimation, produce non-normal confidence intervals, and can
be more robust in small samples. Bootstrap validity still depends on matching
spatial dependence.

\subsection{Choice of resampling scheme}

\paragraph{(i) Site (pairs) bootstrap.}
Resample site indices $\ell\in\{1,\ldots,L\}$ with replacement and re-estimate.
This is credible only if sites are approximately exchangeable after
conditioning on the maintained model.

\paragraph{(ii) Cluster / spatial block bootstrap.}
Partition sites into spatial blocks or clusters and resample clusters with
replacement. This preserves within-cluster dependence and is a reasonable
default in spatial diffusion applications.

\paragraph{(iii) Parametric bootstrap.}
Simulate a full synthetic diffusion history from the fitted model, extract
synthetic arrival times at the observed site locations, and re-estimate. This
preserves the global dependence structure implied by the diffusion model but is
model-based.

\subsection{What distribution does the bootstrap approximate?}

\paragraph{Large-$N$ (oracle) target.}
If $L/N\to 0$ and the resampling scheme is valid for the dependence structure,
then the bootstrap distribution of
$\sqrt{L}(\hat\theta^*_{L,N}-\hat\theta_{L,N})$ approximates the sampling
distribution of $\sqrt{L}(\tilde\theta_L-\theta_0)$.

\paragraph{Finite-$N$ target.}
If $N$ is not large, the bootstrap should mimic the actual reported estimator.
A clean default is to set $N^*=N$ inside the bootstrap. Reducing $N^*$ is a
computational shortcut that changes the estimator whose distribution is being
approximated.

\subsection{Confidence intervals}

Let $\hat\theta^{*(b)}$ denote the estimate from bootstrap replication
$b=1,\ldots,B$. For a scalar component $\hat\theta_k$:

\begin{itemize}
  \item \textbf{Bootstrap standard error:}
  \[
  \widehat{\mathrm{se}}^{\mathrm{boot}}_k
  = \sqrt{\frac{1}{B-1}\sum_{b=1}^B\bigl(\hat\theta_k^{*(b)}-\bar\theta_k^*\bigr)^2},
  \qquad \bar\theta_k^*=\frac{1}{B}\sum_{b=1}^B\hat\theta_k^{*(b)}.
  \]

  \item \textbf{Percentile interval:}
  $[\hat\theta_k^{*(\alpha/2)},\,\hat\theta_k^{*(1-\alpha/2)}]$.

  \item \textbf{Basic interval:}
  $[2\hat\theta_k-\hat\theta_k^{*(1-\alpha/2)},\;2\hat\theta_k-\hat\theta_k^{*(\alpha/2)}]$.

  \item \textbf{Normal interval:}
  $\hat\theta_k \pm z_{1-\alpha/2}\,\widehat{\mathrm{se}}^{\mathrm{boot}}_k$.

  \item \textbf{BCa interval:}
  preferred when available and $B$ is large (software-supported).
\end{itemize}

\begin{remark}[Recommended default]
For reporting, BCa is a strong generic default when available.
When BCa is inconvenient, percentile or basic intervals are reasonable
nonparametric defaults. Normal intervals are useful for compact summary tables
but should not be expected to capture skewness.
\end{remark}

\subsection{Implementation rules}

A robust bootstrap implementation should follow these rules:
\begin{enumerate}
  \item \textbf{Pass site coordinates explicitly.} The simulator should return
        arrival times at a supplied list of coordinates. This avoids hidden
        dependence on a global site ordering and makes repeated sites transparent.

  \item \textbf{Use common random numbers (CRN) within each optimization.}
        Within a given estimation run (original sample or bootstrap replication),
        fix a seed vector $\{s_n\}_{n=1}^{N}$ and reuse it for every evaluation
        of $Q_{L,N}(\theta)$ as $\theta$ varies.

  \item \textbf{Use fresh seed vectors across replications when MC error is part of
        the inferential target.} If the aim is to mimic a finite-$N$ estimator,
        each replication should use its own seed vector. If MC error is negligible,
        reusing the same seed vector across replications is acceptable.

  \item \textbf{Reapply any data-dependent origin rule.} If the procedure selects
        the diffusion origin(s) from the data (e.g. earliest observed site), apply
        the same origin rule inside each replication.

  \item \textbf{Track optimization failures.} Store exit flags, drop failed
        replications, and report the effective number of successful replications
        $B_{\mathrm{eff}}$.
\end{enumerate}

\section{MATLAB implementation templates}

This section provides MATLAB templates implementing:
(i) point estimation with CRN, (ii) analytical sandwich standard errors,
and (iii) cluster/bootstrap inference.

The code assumes you have a function
\texttt{run\_simulator(theta, geo\_layers, origin, site\_coords)}
that returns an $L\times 1$ vector of simulated arrival times at the supplied
coordinates.

\subsection{Seeded simulator wrappers}

\begin{lstlisting}
function Y = run_simulator_seeded(theta, geo_layers, origin, site_coords, seed)
    rng(seed, 'twister');
    Y = run_simulator(theta, geo_layers, origin, site_coords);
end

function Y_hat = simulate_average(theta, geo_layers, origin, site_coords, seeds)
    N = numel(seeds);
    for n = 1:N
        Y_n = run_simulator_seeded(theta, geo_layers, origin, ...
                                   site_coords, seeds(n));
        if n == 1
            L = numel(Y_n);
            Y_sum = zeros(L, 1);
        end
        Y_sum = Y_sum + Y_n;
    end
    Y_hat = Y_sum / N;
end
\end{lstlisting}

\subsection{Objective function with CRN}

\begin{lstlisting}
function mse = objective_crn(theta, Y_obs, geo_layers, origin, site_coords, seeds)
    Y_hat = simulate_average(theta, geo_layers, origin, site_coords, seeds);
    e = Y_obs - Y_hat;
    mse = mean(e.^2);
end
\end{lstlisting}

\subsection{Point estimation}

\begin{lstlisting}
theta0 = [-1.2; 0; 0];
N_opt  = 1000;

rng(24680, 'twister');
seeds_obj = randi(2^31 - 1, N_opt, 1);

opts = optimoptions('fminunc', ...
    'Algorithm',           'quasi-newton', ...
    'Display',             'iter', ...
    'MaxIterations',       500, ...
    'OptimalityTolerance', 1e-6, ...
    'FiniteDifferenceType','central');

obj = @(th) objective_crn(th, Y_obs, geo_layers, origin, ...
                          site_coords, seeds_obj);

[theta_hat, mse_hat, exitflag, output] = fminunc(obj, theta0, opts);
\end{lstlisting}

\subsection{Finite-difference Jacobian with CRN}

Finite differences should reuse the same seed vector for $\theta+h e_k$ and
$\theta-h e_k$.

\begin{lstlisting}
function J = estimate_jacobian_crn(theta_hat, geo_layers, origin, ...
                                   site_coords, seeds_grad, h)
    if nargin < 6 || isempty(h)
        h = 1e-4;
    end

    K = numel(theta_hat);
    Y0 = simulate_average(theta_hat, geo_layers, origin, ...
                          site_coords, seeds_grad);
    L = numel(Y0);
    J = zeros(L, K);

    for k = 1:K
        step = h * max(abs(theta_hat(k)), 1);
        th_p = theta_hat; th_m = theta_hat;
        th_p(k) = th_p(k) + step;
        th_m(k) = th_m(k) - step;

        Yp = simulate_average(th_p, geo_layers, origin, ...
                              site_coords, seeds_grad);
        Ym = simulate_average(th_m, geo_layers, origin, ...
                              site_coords, seeds_grad);

        J(:, k) = (Yp - Ym) / (2 * step);
    end
end
\end{lstlisting}

\subsection{Sandwich variance assembly}

\begin{lstlisting}
function out = compute_variance(theta_hat, Y_obs, geo_layers, origin, ...
                                site_coords, seeds_obj, seeds_grad)
    K = numel(theta_hat);
    L = numel(Y_obs);

    Y_hat = simulate_average(theta_hat, geo_layers, origin, ...
                             site_coords, seeds_obj);
    e = Y_obs - Y_hat;
    J = estimate_jacobian_crn(theta_hat, geo_layers, origin, ...
                              site_coords, seeds_grad);

    Gamma = (J' * J) / L;
    if rcond(Gamma) < 1e-10
        warning('Gamma is ill-conditioned; using pinv.');
        Ginv = pinv(Gamma);
    else
        Ginv = Gamma \ eye(K);
    end

    % Heteroskedastic site-i.i.d. meat
    Meat_iid = (J' * diag(e.^2) * J) / L;
    V_iid = Ginv * Meat_iid * Ginv;      % var of sqrt(L)(theta_hat-theta0)
    se_iid = sqrt(diag(V_iid) / L);      % SE for theta_hat

    % Homoskedastic shortcut
    sigma2 = (e' * e) / max(L - K, 1);
    V_homo = sigma2 * Ginv;             % var of sqrt(L)(theta_hat-theta0)
    se_homo = sqrt(diag(V_homo) / L);

    out.resid = e;
    out.J = J;
    out.Gamma = Gamma;
    out.Meat_iid = Meat_iid;
    out.V_iid = V_iid;
    out.se_iid = se_iid;
    out.sigma2 = sigma2;
    out.V_homo = V_homo;
    out.se_homo = se_homo;
end
\end{lstlisting}

\subsection{Cluster-robust meat and SE}

\begin{lstlisting}
function Meat_cluster = cluster_meat(J, e, cluster_id)
    L = numel(e);
    K = size(J, 2);
    Meat_cluster = zeros(K, K);

    groups = unique(cluster_id(:))';
    for g = groups
        idx = (cluster_id == g);
        s_g = J(idx, :)' * e(idx);
        Meat_cluster = Meat_cluster + s_g * s_g';
    end

    Meat_cluster = Meat_cluster / L;
end
\end{lstlisting}

To convert the cluster meat into cluster-robust standard errors:

\begin{lstlisting}
Meat_cluster = cluster_meat(out_var.J, out_var.resid, cluster_id);
K = numel(theta_hat);
L = numel(Y_obs);
Ginv = out_var.Gamma \ eye(K);
V_cluster = Ginv * Meat_cluster * Ginv;  % var of sqrt(L)(theta_hat-theta0)
se_cluster = sqrt(diag(V_cluster) / L);
\end{lstlisting}

\subsection{Simulation-noise diagnostic}

A simple diagnostic checks whether MC error is small relative to the overall
fit. It estimates the average diagonal variance of $\hat Y_{\ell,n}(\hat\theta)$
across MC draws and reports an ``inflation ratio''.

\begin{lstlisting}
function rho = sim_inflation_ratio(theta_hat, Y_obs, geo_layers, origin, ...
                                   site_coords, seeds_diag)
    L = numel(Y_obs);
    N = numel(seeds_diag);
    Y_draws = zeros(L, N);

    for n = 1:N
        Y_draws(:, n) = run_simulator_seeded(theta_hat, geo_layers, ...
                           origin, site_coords, seeds_diag(n));
    end

    Y_hat = mean(Y_draws, 2);
    total_var = mean((Y_obs - Y_hat).^2);
    sim_var_avgdiag = mean(var(Y_draws, 1, 2)) / N;  % population normalization
    rho = sim_var_avgdiag / total_var;

    fprintf('Simulation inflation ratio (avg diagonal): %.3f\n', rho);
end
\end{lstlisting}

Large values of $\rho$ suggest increasing $N$.

\subsection{Origin-selection helper}

\begin{lstlisting}
function origin_b = choose_origin(coords_b, Y_b, origin_rule, origin_fixed)
    switch lower(origin_rule)
        case 'fixed'
            origin_b = origin_fixed;
        case 'earliest'
            [~, i0] = min(Y_b);
            origin_b = coords_b(i0, :);
        otherwise
            error('Unknown origin_rule. Customize for your application.');
    end
end
\end{lstlisting}

\subsection{Generic cluster bootstrap with CRN}

The routine below implements site bootstrap as the special case
\texttt{cluster\_id = (1:L)'}. For a spatial block bootstrap, set
\texttt{cluster\_id} to geographic block labels.

\begin{lstlisting}
function out = run_cluster_bootstrap(theta_hat, Y_obs, site_coords, ...
                                     cluster_id, geo_layers, ...
                                     origin_rule, origin_fixed, ...
                                     B, N_boot, alpha)
    K = numel(theta_hat);
    theta_boot = NaN(K, B);
    exitflag_boot = NaN(B, 1);

    opts = optimoptions('fminunc', ...
        'Algorithm',           'quasi-newton', ...
        'Display',             'off', ...
        'MaxIterations',       300, ...
        'OptimalityTolerance', 1e-5, ...
        'FiniteDifferenceType','central');

    groups = unique(cluster_id(:), 'stable');
    G = numel(groups);
    cluster_list = cell(G, 1);
    for g = 1:G
        cluster_list{g} = find(cluster_id == groups(g));
    end

    rng(12345, 'twister');
    draw_bank = randi(G, G, B);               % resampling draws
    seed_bank = randi(2^31 - 1, N_boot, B);  % CRN seed vectors

    parfor b = 1:B
        draw_idx = draw_bank(:, b);
        idx = vertcat(cluster_list{draw_idx});

        Y_b = Y_obs(idx);
        coords_b = site_coords(idx, :);
        origin_b = choose_origin(coords_b, Y_b, origin_rule, origin_fixed);
        seeds_b = seed_bank(:, b);

        obj_b = @(th) objective_crn(th, Y_b, geo_layers, origin_b, ...
                                   coords_b, seeds_b);
        [theta_b, ~, exitflag] = fminunc(obj_b, theta_hat, opts);

        if exitflag > 0
            theta_boot(:, b) = theta_b;
        end
        exitflag_boot(b) = exitflag;
    end

    success = all(isfinite(theta_boot), 1);
    B_eff = nnz(success);
    if B_eff < B
        warning(['%d of %d bootstrap replications failed and were dropped.'], ...
                 B - B_eff, B);
    end
    if B_eff < 50
        error('Too few successful bootstrap replications.');
    end

    theta_ok = theta_boot(:, success);
    se_boot = std(theta_ok, 0, 2);
    q_lo = quantile(theta_ok, alpha / 2, 2);
    q_hi = quantile(theta_ok, 1 - alpha / 2, 2);

    ci.pct = [q_lo, q_hi];
    ci.basic = [2 * theta_hat - q_hi, 2 * theta_hat - q_lo];

    z = norminv(1 - alpha / 2);
    ci.norm = [theta_hat - z * se_boot, theta_hat + z * se_boot];

    out.theta_boot = theta_boot;
    out.se_boot = se_boot;
    out.ci = ci;
    out.success = success;
    out.exitflag = exitflag_boot;
    out.B_eff = B_eff;
end
\end{lstlisting}

\subsection{Reporting helper}

\begin{lstlisting}
function report_bootstrap(theta_hat, out, param_names)
    fprintf('\nSuccessful bootstrap replications: %d\n', out.B_eff);
    fprintf('%s\n', repmat('-', 1, 84));
    fprintf('%-18s %10s %10s %20s %20s\n', ...
        'Parameter', 'Estimate', 'SE boot', '95% CI basic', '95% CI pct');
    fprintf('%s\n', repmat('-', 1, 84));

    for k = 1:numel(theta_hat)
        fprintf('%-18s %10.4f %10.4f [%7.4f,%7.4f] [%7.4f,%7.4f]\n', ...
            param_names{k}, theta_hat(k), out.se_boot(k), ...
            out.ci.basic(k, 1), out.ci.basic(k, 2), ...
            out.ci.pct(k, 1), out.ci.pct(k, 2));
    end
    fprintf('%s\n', repmat('-', 1, 84));
end
\end{lstlisting}

\subsection{Bootstrap convergence diagnostic}

A simple diagnostic plots the running bootstrap SE as the number of successful
replications increases.

\begin{lstlisting}
function check_bootstrap_convergence(theta_boot_ok)
    [K, B_eff] = size(theta_boot_ok);
    se_running = NaN(K, B_eff);

    for b = 2:B_eff
        se_running(:, b) = std(theta_boot_ok(:, 1:b), 0, 2);
    end

    figure;
    for k = 1:K
        subplot(K, 1, k);
        plot(2:B_eff, se_running(k, 2:end), 'LineWidth', 1.5);
        xlabel('Successful bootstrap replications');
        ylabel(sprintf('SE(theta_%d)', k));
        title(sprintf('Running bootstrap SE for parameter %d', k));
        grid on;
    end
end
\end{lstlisting}

\subsection{Main script skeleton}

\begin{lstlisting}
%% User choices
param_names = {'theta_0', 'theta_sea', 'theta_prec'};
alpha = 0.05;
B = 999;

N_opt  = 1000;
N_boot = N_opt;  % recommended default: N_boot = N_opt

% For site bootstrap use cluster_id = (1:numel(Y_obs))'.
% For spatial bootstrap replace this by block/cluster labels.
cluster_id = (1:numel(Y_obs))';

% Origin selection
origin_rule  = 'fixed';
origin_fixed = origin;
% origin_rule  = 'earliest';
% origin_fixed = [];

%% Point estimate
rng(24680, 'twister');
seeds_obj = randi(2^31 - 1, N_opt, 1);
origin_base = choose_origin(site_coords, Y_obs, origin_rule, origin_fixed);

obj = @(th) objective_crn(th, Y_obs, geo_layers, origin_base, ...
                          site_coords, seeds_obj);
[theta_hat, mse_hat] = fminunc(obj, theta0, opts);

%% Analytical sandwich SE (optional)
rng(13579, 'twister');
seeds_grad = randi(2^31 - 1, 3000, 1);
out_var = compute_variance(theta_hat, Y_obs, geo_layers, origin_base, ...
                           site_coords, seeds_obj, seeds_grad);

%% Optional simulation-noise diagnostic
rng(54321, 'twister');
seeds_diag = randi(2^31 - 1, 1000, 1);
rho = sim_inflation_ratio(theta_hat, Y_obs, geo_layers, origin_base, ...
                          site_coords, seeds_diag);

%% Bootstrap
out_boot = run_cluster_bootstrap(theta_hat, Y_obs, site_coords, ...
                                 cluster_id, geo_layers, ...
                                 origin_rule, origin_fixed, ...
                                 B, N_boot, alpha);

theta_ok = out_boot.theta_boot(:, out_boot.success);
check_bootstrap_convergence(theta_ok);
report_bootstrap(theta_hat, out_boot, param_names);
\end{lstlisting}

\section{Reporting guidance}

\begin{itemize}
  \item \textbf{Choose $N$ to control MC error.} If inference is intended to reflect
        uncertainty from the data, increase $N$ until MC variation in
        $\bar Y_{\ell,N}(\hat\theta)$ is small relative to residual variation.
        The inflation diagnostic above is one practical check.

  \item \textbf{Prefer cluster/block procedures when spatial dependence is plausible.}
        In diffusion applications, this is often more credible than treating sites
        as independent.

  \item \textbf{Use $B\ge 999$ for final nonparametric intervals.} Smaller $B$ can be
        used for debugging.

  \item \textbf{Transform parameters by transforming draws.} If reporting a transformed
        object (e.g., an implied speed), apply the transformation to each bootstrap
        draw and then compute quantiles.
\end{itemize}

\section{Summary}

Analytical (sandwich) inference and bootstrap inference are both natural for
simulation-based diffusion estimators.

\begin{itemize}
  \item The analytical approach is compact and can be fast once a stable Jacobian
        estimate is available. It requires careful attention to (i) Jacobian
        computation (finite differences with CRN), (ii) the dependence structure
        across sites, and (iii) the role of simulation size $N$.

  \item The bootstrap approach avoids explicit Jacobians and supports non-normal
        intervals, but still requires choosing a resampling scheme that matches
        spatial dependence (site vs block/cluster) and handling simulation seeds
        carefully.
\end{itemize}

In both approaches, using common random numbers within each optimization run is
crucial for numerical stability.

\end{document}
